% !TEX program = lualatex
\documentclass[a4paper,11pt]{article}
\usepackage{amsmath,amssymb,amsthm,mathtools}
\usepackage{geometry}
\usepackage{enumitem}
\usepackage{hyperref}
\geometry{margin=25mm}

\newtheorem{definition}[section]{定義}
\newtheorem{proposition}[definition]{命題}
\newtheorem{theorem}[definition]{定理}
\newtheorem{example}[definition]{例}
\newtheorem{counterexample}[definition]{反例}
\newtheorem{remark}[definition]{注意}

\newcommand{\ECA}{\mathrm{ECA}}
\newcommand{\Ax}{\mathrm{Ax}}
\newcommand{\Solver}{\mathsf{Solver}}
\newcommand{\Th}{\mathrm{Th}}
\newcommand{\Pow}{\mathcal{P}}
\newcommand{\ZFC}{\mathrm{ZFC}}

% ECA_18～24
\title{ECAが連続体濃度を扱うことが可能か検証するための例・反例・検証条件}
\author{}
\date{}

\begin{document}
\maketitle

\section{目的}

本稿では、ECA（拡張連続性公理）の既存構造そのものが、連続体濃度
\[
2^{\aleph_0}
\]
に到達するだけの濃度をECAの数学的構造を経て導出し得るか検証し、ECAが既に持つ構造だけから連続体濃度が導出されるか否かを、例と反例を用いて可能な限り形にすることを目的とする。


対象は、これまで整理した
\[
\mathfrak U_{\mathrm{ECA}}
\longrightarrow
\operatorname{Ax}(U)
\longrightarrow
\mathcal{R}^*
\longrightarrow
\mathcal{P}(\mathcal{R}^*)
\longrightarrow
I
\longrightarrow
\Solver_I
\longrightarrow
\mathcal{Z}
\]
という構造が対象となる。

各段階について、既存のECA定義から何が導けるか、何がまだ導けないか、どのような例がその構造を実現するか、どのような反例によって過剰な推測を排除できるかを確認する。

\section{現在の基本構造}

ECAでは、数学的宇宙の全体を
\[
\mathfrak U_{\mathrm{ECA}}
\]
として扱い、その要素を
\[
U\in\mathfrak U_{\mathrm{ECA}}
\]
とする。

各数学的宇宙に対応する公理候補集合を
\[
\operatorname{Ax}(U)
\]
とし、ECA全体で扱う公理および公理スキームの候補集合を
\[
\Ax
=
\bigcup_{U\in\mathfrak U_{\mathrm{ECA}}}\operatorname{Ax}(U)
\]
としている。

ただし、この式を文字通り集合の集まりとして扱うためには、対象となる数学的宇宙や候補集合の集合論的な型を適切に考える必要がある。

公理選択については
\[
I\subseteq\mathcal{P}(\Ax)
\]
とし、
\[
i\in I
\]
をECAにおける許容された公理集合とする。

さらに
\[
B(i)\subseteq i\subseteq\Ax
\]
を維持する。

Solverは
\[
s=(\mathcal{A}_s,\mathcal{R}_s)
\]
とし、
\[
\Solver_I
=
\left\{
s=(\mathcal{A}_s,\mathcal{R}_s)
\mid
\mathcal{A}_s\in I
\right\}
\]
とする。

数学的構造写像は
\[
\mathcal{T}:\Solver_I\to\Th
\]
であり、Solver $s$ が表現する数学的構造を
\[
\mathcal{T}(s)
\]
とする。

ZFCを表現するSolver集合は
\[
\boxed{
\mathcal{Z}
=
\left\{
s\in\Solver_I
\mid
\mathcal{T}(s)\cong_{\mathrm{Th}}\ZFC
\right\}
}
\]
である。

したがって、
\[
|\Ax|,\quad |I|,\quad |\Solver_I|,\quad |\mathcal{Z}|
\]
は同一の濃度ではなく、それぞれ別個に追跡する必要がある。

\section{基本原理}

連続体濃度をECAの構造から得る典型的な経路は、可算無限公理候補族
\[
\mathcal{R}^*
=
\{r_0,r_1,r_2,\ldots\}
\]
を取り出し、
\[
|\mathcal{R}^*|=\aleph_0
\]
から
\[
|\mathcal{P}(\mathcal{R}^*)|
=
2^{\aleph_0}
\]
を利用することである。

しかし、
\[
|\mathcal{R}^*|=\aleph_0
\]
から直ちに
\[
|\mathcal{Z}|\ge2^{\aleph_0}
\]
とはならない。

少なくとも次の段階を区別しなければならない。
\[
\operatorname{Ax}(U)
\longrightarrow
\mathcal{R}^*
\longrightarrow
\mathcal{P}(\mathcal{R}^*)
\longrightarrow
I
\longrightarrow
\Solver_I
\longrightarrow
\mathcal{Z}.
\]

特に、
\[
|\operatorname{Ax}(U)|\ge\aleph_0
\]
は公理候補の存在に関する主張であり、
\[
|I|\ge2^{\aleph_0}
\]
は選択族の濃度に関する主張であり、
\[
|\mathcal{Z}|\ge2^{\aleph_0}
\]
はZFC表現Solverの濃度に関する主張である。

これらを同一視してはならない。

\section{例1：可算な公理候補族}

まず、
\[
\exists U\in\mathfrak U_{\mathrm{ECA}},
\qquad
\exists\mathcal{R}^*\subseteq\operatorname{Ax}(U)
\]
について
\[
|\mathcal{R}^*|=\aleph_0
\]
となる例を考える。

ここで重要なのは、公理スキームをどのような対象として
\[
\operatorname{Ax}(U)
\]
に追加するかである。

公理スキームそのものを1個の要素として扱う場合と、そのスキームから得られる個々のインスタンスを公理候補として扱う場合では、集合としての
\[
\operatorname{Ax}(U)
\]
の構造が異なる。

したがって、ECAの「公理および公理スキーム」という定義だけから、
\[
|\operatorname{Ax}(U)|\ge\aleph_0
\]
を自然に導くことはできない。

一方、ある公理スキームの異なるインスタンスを個別の候補として扱う形式化を採用できれば、
\[
\mathcal{R}^*
=
\{r_n\mid n\in\mathbb{N}\}
\subseteq\operatorname{Ax}(U)
\]
という可算族を構成できる可能性がある。

ここでは、その形式化がECAの既存定義から自然に導かれるかを検証対象とする。

\section{例2：可算族の冪集合}

可算無限族
\[
\mathcal{R}^*=\{r_n\mid n\in\mathbb{N}\}
\]
が得られたと仮定する。

このときカントールの定理から
\[
|\mathcal{P}(\mathcal{R}^*)|
=
2^{|\mathcal{R}^*|}
=
2^{\aleph_0}
\]
であることがわかる。

したがって、
\[
A\subseteq\mathcal{R}^*
\]
ごとに異なる対象をECAで追加することができれば、連続体濃度を持つ族を導出するための数学的基礎は存在する。

しかし、
\[
\mathcal{P}(\mathcal{R}^*)
\]
の存在と、
\[
\mathcal{P}(\mathcal{R}^*)\subseteq I
\]
は別の事実である。

\section{例3：選択族への埋め込み}

\[
\mathcal{A}_{\mathrm{ZFC}}\subseteq\Ax,
\qquad
\mathcal{R}^*\cap\mathcal{A}_{\mathrm{ZFC}}=\varnothing
\]
とする。

各
\[
A\subseteq\mathcal{R}^*
\]
について
\[
i_A
=
\mathcal{A}_{\mathrm{ZFC}}\cup A
\]
を定める。

ここで
\[
\boxed{
\forall A\subseteq\mathcal{R}^*,
\qquad
i_A\in I
}
\]
が成立すれば、
\[
\Gamma:
\mathcal{P}(\mathcal{R}^*)\to I,
\qquad
\Gamma(A)=i_A
\]
を定義できる。

\[
\mathcal{R}^*\cap\mathcal{A}_{\mathrm{ZFC}}=\varnothing
\]
なら
\[
A\neq B
\Longrightarrow
i_A\neq i_B
\]
であるから、
\[
\Gamma
\]
は単射となる。

したがって、
\[
\boxed{
|I|\ge2^{\aleph_0}
}
\]
が得られる。

\section{反例1：$I$ が小さい場合}

\begin{counterexample}
\[
\operatorname{Ax}(U)
\]
に可算無限集合
\[
\mathcal{R}^*
\]
が存在しても、
\[
I=\{\mathcal{A}_{\mathrm{ZFC}}\}
\]
とすることは、現在の
\[
I\subseteq\mathcal{P}(\Ax)
\]
という定義だけからは排除されない。

この場合、
\[
\mathcal{P}(\mathcal{R}^*)
\]
は確かに連続体濃度を持つが、選択可能な公理集合の族は増えない。

したがって、
\[
|\operatorname{Ax}(U)|\ge\aleph_0
\]
だけから
\[
|I|\ge2^{\aleph_0}
\]
を導くことはできない。
\end{counterexample}

\section{反例2：$I$ が大きくても $\mathcal{Z}$ が小さい場合}

\begin{counterexample}
\[
|I|\ge2^{\aleph_0}
\]
であっても、ZFCを表現するSolverに対応する選択が1つしかない場合、
\[
|\mathcal{Z}|=1
\]
となり得る。

すなわち、
\[
|I|\ge2^{\aleph_0}
\not\Rightarrow
|\mathcal{Z}|\ge2^{\aleph_0}.
\]
\end{counterexample}

したがって、\(\mathcal{Z}\) の濃度を考慮するには、選択族の大きさだけでなく、
\[
\mathcal{T}(s)\cong_{\mathrm{Th}}\ZFC
\]
を維持したまま異なるSolverを構成できるかを確認する必要がある。

\section{例4：ZFC保存族}

可算無限族
\[
\mathcal{R}^*
=
\{r_n\mid n\in\mathbb{N}\}
\]

と、Solver内の固定された推論規則・構成制約

\[
\mathcal{R}_0
\]
を考える。

各
\[
A\subseteq\mathcal{R}^*
\]
について
\[
s_A
=
\left(
\mathcal{A}_{\mathrm{ZFC}}\cup A,
\mathcal{R}_0
\right)
\]
を定める。

このとき、次の条件を考慮する。
\begin{align}
&\forall A\subseteq\mathcal{R}^*,
\quad
\mathcal{A}_{\mathrm{ZFC}}\cup A\in I,
\tag{S2}\\
&\forall A\subseteq\mathcal{R}^*,
\quad
\mathcal{T}(s_A)
\cong_{\mathrm{Th}}\ZFC.
\tag{S3}
\end{align}

(S2) は選択可能性、(S3) はZFC保存性を表す。

さらに
\[
\mathcal{R}^*\cap\mathcal{A}_{\mathrm{ZFC}}
=
\varnothing
\]
なら、
\[
A\neq B
\Longrightarrow
s_A\neq s_B
\]
である。

したがって、
\[
\Phi:
\mathcal{P}(\mathcal{R}^*)\to\mathcal{Z},
\qquad
\Phi(A)=s_A
\]
は単射となる。

よって、これらの条件が実際に構成できれば、
\[
\boxed{
|\mathcal{Z}|
\ge
2^{\aleph_0}
}
\]
が得られる。

ここでも、この不等式は現段階でECAから自然に得られた定理ではなく、(S2)、(S3)等の成立を確認した場合の帰結である。

\section{検証項目：ZFC保存性の意味}

(S3) を成立させるには、各
\[
r\in\mathcal{R}^*
\]
が単独で「ZFCを壊さない」だけでは足りない。

必要なのは、
\[
\boxed{
\forall A\subseteq\mathcal{R}^*,
\quad
\mathcal{T}
\left(
\mathcal{A}_{\mathrm{ZFC}}\cup A,\mathcal{R}_0
\right)
\cong_{\mathrm{Th}}
\ZFC
}
\]
である。

つまり、\(\mathcal{R}^*\) の各要素をどのように選んでも、選択された全体がZFCと同型な数学的構造を表現し続ける必要がある。

この検証では、少なくとも以下を分離する必要がある。
\begin{enumerate}
  \item ZFCの定理であること。
  \item ZFCの公理として既に採用されていること。
  \item ECAの$\Ax$に公理候補として含まれていること。
  \item その候補を追加しても数学的構造が保存されること。
  \item 任意の部分集合を同時に追加しても数学的構造が保存されること。
\end{enumerate}

特に
\[
\ZFC\vdash r
\]
から直ちに
\[
r\in\Ax
\]
とはならない。

したがって、「ZFCの定理を大量に持ってくればよい」という議論だけでは、ECA内部の
\[
\mathcal{R}^*\subseteq\Ax
\]
を構成したことにはならない。

\section{例5：Solver構造側からの濃度の導出}

公理選択を変えずに、
\[
\mathcal{A}_s=\mathcal{A}_t
\]
でありながら
\[
\mathcal{R}_s\neq\mathcal{R}_t
\]
となるSolverを考えることもできる。

例えば
\[
s_R=(\mathcal{A}_0,\mathcal{R}_R)
\]
という族を構成し、
\[
R\neq R'
\Longrightarrow
\mathcal{R}_R\neq\mathcal{R}_{R'}
\]
かつ
\[
\mathcal{T}(s_R)
\cong_{\mathrm{Th}}
\ZFC
\]
を成立させられれば、Solverそのものの濃度を増加させることができる。

この場合、\(\mathfrak A\) の像は増加しない可能性がある。

したがって、
\[
\left|\operatorname{Im}
(\mathfrak A|_{\mathcal{Z}})\right|
\]
と
\[
|\mathcal{Z}|
\]
を分離して追跡する必要がある。

\section{検証項目：$\mathcal{Z}$ のファイバー分解}

\[
\mathcal{Z}_A
=
\left\{
s\in\mathcal{Z}
\mid
\mathfrak A(s)=A
\right\}
\]
と定義する。

すると
\[
\mathcal{Z}
=
\bigcup_{A\in
\operatorname{Im}(\mathfrak A|_{\mathcal{Z}})}
\mathcal{Z}_A.
\]

ここで、
\[
\left|\operatorname{Im}
(\mathfrak A|_{\mathcal{Z}})\right|
\]
は異なる公理選択の数を表し、
\[
|\mathcal{Z}_A|
\]
は1つの公理選択に対して存在する異なるSolver構造の数を表す。

したがって、
\[
|\mathcal{Z}|\ge2^{\aleph_0}
\]
となる経路は、少なくとも概念的には、
\[
\left|\operatorname{Im}
(\mathfrak A|_{\mathcal{Z}})\right|
\ge2^{\aleph_0}
\]
による場合と、
\[
\exists A\quad
|\mathcal{Z}_A|\ge2^{\aleph_0}
\]
による場合とに分けて考えることができる。

さらに両方が組み合わさる可能性もある。

\section{反例3：同じ公理選択にすべてが潰れる場合}

\begin{counterexample}
異なる候補集合
\[
A\neq B
\]
を用意しても、それらを同じSolverとして構成できるような構造を採用すれば、
\[
s_A=s_B
\]
となる可能性がある。

この場合、
\[
\mathcal{P}(\mathcal{R}^*)
\]
から\(\mathcal{Z}\)への単射は成立しない。

したがって、連続体濃度を得るには、単に異なる選択を列挙するだけでなく、
\[
A\neq B
\Longrightarrow
s_A\neq s_B
\]
を保証する必要がある。

ECA_6で定義したSolverの構造としては、
\[
s=(\mathcal{A}_s,\mathcal{R}_s)
\]
であるため、特に
\[
\mathcal{A}_{s_A}\neq\mathcal{A}_{s_B}
\]
または
\[
\mathcal{R}_{s_A}\neq\mathcal{R}_{s_B}
\]
を利用してSolverの非同一性を確認できる。
\end{counterexample}

\section{上界の考慮}
\(\mathcal{Z}\subseteq\Solver_I\) であるため、
\[
\boxed{
|\mathcal{Z}|
\le
|\Solver_I|
}
\]
は直ちに得られる。

一方、
\[
|\Solver_I|\le2^{\aleph_0}
\]
を現在のECAから導く根拠はまだ確立されていない。

したがって、連続体濃度の下界を構成できたとしても、
\[
|\mathcal{Z}|=2^{\aleph_0}
\]
とは直ちには言えない。

現段階での一般的な範囲は、条件付きで
\[
\boxed{
2^{\aleph_0}
\le
|\mathcal{Z}|
\le
|\Solver_I|
}
\]
という形に留める。

上界をさらに狭めるには、
\[
|\Ax|,\quad |I|,\quad |\Solver_I|
\]
の濃度を別途定義、検証する必要がある。

\section{検証表}

今後の検証は次の順序で行う。
ただし、下記で列挙されている検証項目はこの文書、およびECAの文書や他の文書で定義した内容、判明したことに直接影響を与えるものではなく、同時に下記の検証項目の立証の可否によっては変わる可能性があることに留意する。
また、この検証項目はあくまでも「可能性の話」として捉え、下記の検証項目の立証の可否がないことを理由としてECAの他の文書で既に定義されたものや、その数学的解釈などの是非を問うことはしないこと。

\begin{enumerate}
  \item \textbf{数学的宇宙の検証}
  \[
  U\in\mathfrak U_{\mathrm{ECA}}
  \]
  が適切に形式化されているか。

  \item \textbf{公理候補集合の検証}
  \[
  \operatorname{Ax}(U)
  \]
  が集合として定義され、可算族を取り出せるか。

  \item \textbf{可算族の検証}
  \[
  \mathcal{R}^*\subseteq\operatorname{Ax}(U),
  \qquad
  |\mathcal{R}^*|=\aleph_0
  \]
  が構成できるか。

  \item \textbf{冪集合の検証}
  \[
  |\mathcal{P}(\mathcal{R}^*)|=2^{\aleph_0}
  \]
  をECAの後段へ接続できるか。

  \item \textbf{選択可能性の検証}
  \[
  \forall A\subseteq\mathcal{R}^*,
  \quad
  \mathcal{A}_{\mathrm{ZFC}}\cup A\in I
  \]
  を既存構造から得られるか。

  \item \textbf{Solver構成の検証}
  \[
  s_A=(\mathcal{A}_{\mathrm{ZFC}}\cup A,\mathcal{R}_0)
  \]
  を各$A$について構成できるか。

  \item \textbf{ZFC保存性の検証}
  \[
  \forall A\subseteq\mathcal{R}^*,
  \quad
  \mathcal{T}(s_A)\cong_{\mathrm{Th}}\ZFC
  \]
  が成立するか。

  \item \textbf{単射性の検証}
  \[
  A\neq B\Longrightarrow s_A\neq s_B
  \]
  を確認できるか。

  \item \textbf{上界の検証}
  \[
  |\Solver_I|
  \]
  を評価し、
  \[
  |\mathcal{Z}|
  \]
  の上界をどこまで狭められるかを調べる。
\end{enumerate}

\section{現段階で確認できることと未確認事項}

現在のECA構造から確実に分離して述べられるのは、次の点である。
\begin{enumerate}
  \item 可算無限集合\(\mathcal{R}^*\)が存在すれば、
  \[
  |\mathcal{P}(\mathcal{R}^*)|=2^{\aleph_0}
  \]
  である。

  \item しかし、現在の
  \[
  I\subseteq\mathcal{P}(\Ax)
  \]
  だけでは、すべての部分集合が選択可能であるとは言えない。

  \item さらに、\(I\) の濃度が大きくても、ZFCを表現するSolverの濃度が大きいとは限らない。

  \item \(\mathcal{Z}\) の連続体濃度以上の下界には、
  \[
  \mathcal{P}(\mathcal{R}^*)\hookrightarrow\mathcal{Z}
  \]
  のような明示的な単射が必要である。

  \item \(\mathcal{Z}\) の上界については、
  \[
  |\mathcal{Z}|\le|\Solver_I|
  \]
  が得られるが、\(|\Solver_I|\) 自体の上界はまだ固定されていない。
\end{enumerate}

したがって、現時点では
\[
|\mathcal{Z}|\ge2^{\aleph_0}
\]
をECAの確定した定理として記述することはできない。

一方で、これまでの構造化によって、この不等式を得るために必要な条件はかなり明確になっている。

\section{今後の検証方針}

今後は、\(\mathcal{R}^*\) の存在を外部から仮定するのではなく、まず
\[
\operatorname{Ax}(U)
\]
の形式化から検証する。

特に、
\[
\boxed{
\text{公理スキーム}
\longrightarrow
\text{個別の公理候補}
}
\]
という対応がECAの構造内で自然に定義できるかを確認する。

その上で、
\[
\operatorname{Ax}(U)
\supseteq
\mathcal{R}^*
\]
を構成し、
\[
\mathcal{P}(\mathcal{R}^*)
\longrightarrow
I
\longrightarrow
\Solver_I
\longrightarrow
\mathcal{Z}
\]
という順で検証する。

また、各検証項目については、それぞれを「独立した命題」として扱い、それぞれ固有の検証条件を持つものとする。
\[
\boxed{
\begin{array}{c}
|\mathcal{R}^*|=\aleph_0\\[2mm]
\Downarrow\\
|\mathcal{P}(\mathcal{R}^*)|=2^{\aleph_0}\\[2mm]
\Downarrow\quad\text{（選択可能性が必要）}\\
|I|\ge2^{\aleph_0}\\[2mm]
\Downarrow\quad\text{（ZFC保存性とSolver非同一性が必要）}\\
|\mathcal{Z}|\ge2^{\aleph_0}
\end{array}
}
\]




\section{ECA_19 公理スキームから可算無限公理候補族 $\mathcal{R}^*$ を構成する方法}


\section{目的}

ECA_18では、ECAが連続体濃度
\[
2^{\aleph_0}
\]
を導出できるかを検証するため、
\[
\mathfrak U_{\mathrm{ECA}}
\to
\operatorname{Ax}(U)
\to
\mathcal{R}^*
\to
\mathcal P(\mathcal{R}^*)
\to
I
\to
\Solver_I
\to
\mathcal Z
\]
という検証手順を順を追って整理した。

本節では、その最初の重要な段階、すなわち
\[
\boxed{
\operatorname{Ax}(U)
\text{ から }
|\mathcal{R}^*|=\aleph_0
\text{ なる }
\mathcal{R}^*
\text{ を構成できるか}
}
\]
を検証する。

ここでは、最初から
\[
\mathcal{R}^*\subseteq\Ax
\]
を仮定するのではなく、ECA_13からECA_15で導入した
\[
\operatorname{Ax}(U)
\]
と「公理および公理スキーム」という構造から、可算無限候補族を自然に取り出せるかを調べる。

\section{ECAにおける公理候補の位置}

ECAでは数学的宇宙を
\[
U\in\mathfrak U_{\mathrm{ECA}}
\]
とし、それに対応する公理候補集合を
\[
\operatorname{Ax}(U)
\]
とする。

また、ECA全体で扱う候補集合を
\[
\Ax
=
\bigcup_{U\in\mathfrak U_{\mathrm{ECA}}}
\operatorname{Ax}(U)
\]
としている。

このとき
\[
\operatorname{Ax}(U)\subseteq\Ax
\]
である。

したがって、ある
\[
\mathcal{R}^*\subseteq\operatorname{Ax}(U)
\]
を構成できれば、
\[
\mathcal{R}^*\subseteq\Ax
\]
も得られる。

しかし、ここで重要なのは、
\[
\operatorname{Ax}(U)
\]
が単に「Uの内部に存在する公理」を意味するものではないことである。

\[
\operatorname{Ax}(U)
\]
はECAがその数学的宇宙に関連して公理的対象として扱う候補の集合であり、その構成方法自体を形式化する必要がある。

\section{公理と公理スキームの区別}

「公理および公理スキーム」を
\[
\operatorname{Ax}(U)
\]
の要素として扱う場合、少なくとも二つの形式化が考えられる。

\subsection{スキームを一つの対象とする場合}

公理スキーム
\[
\Sigma
\]
を一個の対象として
\[
\Sigma\in\operatorname{Ax}(U)
\]
とする。

この場合、スキームが無限個のインスタンスを生成するとしても、
\[
\Sigma
\]
自体は一個の要素である。

したがって、
\[
\Sigma
\]
の存在だけから
\[
|\operatorname{Ax}(U)|\ge\aleph_0
\]
は導けない。

\subsection{インスタンスを個別候補とする場合}

一方、スキーム
\[
\Sigma(x)
\]
から得られる各インスタンスを
\[
\Sigma(\varphi_n)
\]
のように個別の公理候補として扱う形式化を採用するなら、
\[
\mathcal{R}^*
=
\{
\Sigma(\varphi_n)\mid n\in\mathbb{N}
\}
\]
を構成できる可能性がある。

ただし、そのためには
\[
\Sigma(\varphi_n)\in\operatorname{Ax}(U)
\]
がすべての
\[
n\in\mathbb{N}
\]
について成立することをECAの定義として保証する必要がある。

したがって、

\[
\boxed{
\text{「公理スキームが存在する」}
\not\Rightarrow
\text{「その全インスタンスが }\operatorname{Ax}(U)\text{ の要素である」}
}
\]

という区別と関係をここでは維持する。

\section{構成候補：可算なインスタンス族}

公理スキーム
\[
\Sigma(\varphi)
\]
があり、
\[
\{\varphi_n\mid n\in\mathbb{N}\}
\]
という可算な系列の入力に対して異なるインスタンスを生成するとする。

各
\[
n\in\mathbb{N}
\]
について
\[
r_n=\Sigma(\varphi_n)
\]
と定める。

さらに、
\[
r_n\in\operatorname{Ax}(U)
\]
をすべての
\[
n\in\mathbb{N}
\]
について保証できるなら、
\[
\mathcal{R}^*
=
\{r_n\mid n\in\mathbb{N}\}
\subseteq\operatorname{Ax}(U)
\]
となる。

ここでさらに
\[
n\neq m\Longrightarrow r_n\neq r_m
\]
が成立すれば、
\[
|\mathcal{R}^*|=\aleph_0
\]
である。

したがって、可算無限候補族を得るための最小条件の一つは、

\[
\boxed{
\begin{aligned}
&\forall n\in\mathbb{N},\quad
r_n\in\operatorname{Ax}(U),\\
&\forall n,m\in\mathbb{N},\quad
n\neq m\Rightarrow r_n\neq r_m.
\end{aligned}
}
\]

である。

\section{注意：可算な入力系列だけでは十分でない}

\begin{counterexample}
可算な系列
\[
\varphi_0,\varphi_1,\ldots
\]
が存在しても、
\[
\Sigma(\varphi_n)=\Sigma(\varphi_m)
\]
となる場合には、生成される候補は有限個または一個になり得る。

したがって、
\[
|\{\varphi_n\mid n\in\mathbb{N}\}|=\aleph_0
\]
だけでは
\[
|\mathcal{R}^*|=\aleph_0
\]
を保証しない。
\end{counterexample}

必要なのは、インスタンスそのものの非同一性である。

\section{検証：ZFCにおける公理スキームをどう扱うか}

ECAの
\[
\operatorname{Ax}(U)
\]
を具体的に検証する際には、ZFCの公理スキームを一つの標準例として利用できる。

ただし、これはZFCそのものをECAの公理候補宇宙と同一視することを意味しない。

ここで検証すべきなのは、ECAが

\[
\text{ZFCにおける公理スキーム}
\longrightarrow
\text{個別インスタンス}
\longrightarrow
\operatorname{Ax}(U)
\]

という対応を採用できるかどうかである。

例えば、あるスキーム
\[
\Sigma(\varphi)
\]
について、可算な異なる入力
\[
\varphi_0,\varphi_1,\ldots
\]
を選び、
\[
r_n=\Sigma(\varphi_n)
\]
とする。

このとき、ECA側で
\[
r_n\in\operatorname{Ax}(U)
\]
を保証する仕組みがあれば、可算族
\[
\mathcal{R}^*
\]
の候補となる。

この段階では、それらをZFCの「冗長な公理」として扱う必要はない。

まず必要なのは、あくまで
\[
\operatorname{Ax}(U)
\]
の要素として可算無限族を形成できるかというものである。

\section{別の構成候補：構文的に異なる公理候補}

公理スキームのインスタンス化以外にも、構文的に異なる公理候補を用いる方法がある。

例えば、
\[
r_0,r_1,r_2,\ldots
\]
がそれぞれ異なる有限の構文表現として
\[
\operatorname{Ax}(U)
\]
に登録されるなら、
\[
\mathcal{R}^*
=
\{r_n\mid n\in\mathbb{N}\}
\]
を構成できる。

この方法では、
\[
r_n\neq r_m
\]
という構文上の非同一性を確認するだけでよい場合がある。

ただし、後段でZFC保存性を要求する場合は
\[
\mathcal{T}
\left(
\mathcal{A}_{\mathrm{ZFC}}\cup A,\mathcal{R}_0
\right)
\cong_{\mathrm{Th}}\ZFC
\]
を任意の
\[
A\subseteq\mathcal{R}^*
\]
について検証しなければならない。

したがって、この方法は現段階では「候補族を作る方法」の一つとして扱い、ZFC保存性とは分けて考えるものとする。

\section{反例：ZFCの定理を列挙するだけでは足りない}

\begin{counterexample}
ZFCで証明可能な命題を
\[
\{q_0,q_1,q_2,\ldots\}
\]
と列挙できたとしても、
\[
q_n\in\operatorname{Ax}(U)
\]
が保証されるとは限らない。

\[
\ZFC\vdash q_n
\]
は「q_nがZFCから導出可能である」ことを示すだけであり、
\[
q_n
\]
がECAの公理候補集合に含まれていることを意味しない。

したがって、
\[
\ZFC\vdash q_n
\]
から
\[
\mathcal{R}^*
=
\{q_n\mid n\in\mathbb{N}\}
\subseteq\operatorname{Ax}(U)
\]
を直接結論することはできない。
\end{counterexample}

これははECA_11で確認した
\[
\text{「ZFCの定理」}
\neq
\text{「ECAの公理候補」}
\]
という関係を維持するものである。

\section{検証条件の整理}
現段階で、ECAの構造から
\[
|\mathcal{R}^*|=\aleph_0
\]
を導くために必要な条件は、少なくとも次のように分解できる。

\subsection*{条件 R1：対象となる数学的宇宙}

\[
\exists U\in\mathfrak U_{\mathrm{ECA}}
\]

\subsection*{条件 R2：候補集合の存在}

\[
\operatorname{Ax}(U)
\]

が集合として適切に定義されていること。

\subsection*{条件 R3：可算系列の存在}

\[
\exists\{r_n\}_{n\in\mathbb{N}}
\]

\subsection*{条件 R4：候補集合への所属}

\[
\forall n\in\mathbb{N},
\quad
r_n\in\operatorname{Ax}(U)
\]

\subsection*{条件 R5：相互非同一性}

\[
\forall n,m\in\mathbb{N},
\quad
n\neq m\Rightarrow r_n\neq r_m
\]

R1からR5が成立すれば

\[
\mathcal{R}^* = \{r_n\mid n\in\mathbb{N}\}
\]

について
\[
\boxed{\mathcal{R}^*\subseteq\operatorname{Ax}(U), \qquad |\mathcal{R}^*|=\aleph_0}
\]
が得られる。

\section{ECAから自然に得られる条件か}

現在のECA定義だけから、

\[
\exists U,\quad
|\operatorname{Ax}(U)|\ge\aleph_0
\]
が確定しているとはまだ言えない。

また、
\[
\text{「公理スキームを扱う」}
\]
という記述だけから、
\[
\forall n,\quad r_n\in\operatorname{Ax}(U)
\]
を確定することもできない。

ただし $\operatorname{Ax}(U)$ の形式化をさらに具体化することでR1からR5が導けるか調べる、より厳密には「公理スキームとそのインスタンスの関係をECAの基本構造として明示する」ことができれば、R3からR5が既存構造から導出可能になる可能性がある。

\section{連続体濃度への接続はまだ行わない}

R1からR5によって
\[
|\mathcal{R}^*|=\aleph_0
\]
が得られた。

次の段階では、
\[
|\mathcal P(\mathcal{R}^*)|
=
2^{\aleph_0}
\]
という集合論的事実に加えて、
\[
\mathcal P(\mathcal{R}^*)\to I
\]
という選択構造が必要である。

さらに、
\[
I\to\Solver_I\to\mathcal Z
\]
では、ZFC保存性とSolverの非同一性を検証する必要がある。

したがって、ここまでの成果はあくまでも「$\operatorname{Ax}(U)$ から可算無限候補族を取り出せる条件の明確化」であることに留意する。

\section{次の検証への接続}

R1からR5を検証した後には、
\[
\mathcal{R}^*
\]
の全部分集合
\[
A\subseteq\mathcal{R}^*
\]
をECAの選択構造
\[
I\subseteq\mathcal P(\Ax)
\]
へ接続できるかを調べる。

すなわち、
\[
\mathcal{A}_{\mathrm{ZFC}}\cup A\in I
\]
をどこまで導けるかを検証する。

その後に、
\[
s_A
=
(\mathcal{A}_{\mathrm{ZFC}}\cup A,\mathcal{R}_0)
\]
を構成し、
\[
\mathcal{T}(s_A)\cong_{\mathrm{Th}}\ZFC
\]
を任意の
\[
A\subseteq\mathcal{R}^*
\]
について確認する。

この順序によって、

\[
\boxed{
\operatorname{Ax}(U)
\to
\mathcal{R}^*
\to
\mathcal P(\mathcal{R}^*)
\to
I
\to
\Solver_I
\to
\mathcal Z
}
\]

の各部分を混同せずに検証できる。





\section{ECA_20 $\mathcal{R}^*$ の選択可能性と ECA 構造からの導出条件}


\section{目的}
ECA_19 では、数学的宇宙 $U$ に対応する
$\operatorname{Ax}(U)$ から可算無限公理候補族
\[
\mathcal{R}^*=\{r_n\mid n\in\mathbb{N}\},
\qquad |\mathcal{R}^*|=\aleph_0
\]
を構成できる条件を検討した。

本節では次の段階として、$\mathcal{R}^*$ の全部分集合を ECA の公理選択空間 $I$ において選択可能にできるかを検討する。

現在の
\[
I\subseteq\mathcal{P}(\Ax)
\]
から
\[
\mathcal{P}(\mathcal{R}^*)\subseteq I
\]
を導くことはできない。そこで本節では、強すぎる大域条件を置かず、
\[
\mathcal{R}^*\longrightarrow\mathcal{P}(\mathcal{R}^*)
\longrightarrow I
\]
を検証する。

\section{前提}
ECA の数学的宇宙の集合を $\mathfrak U_{\mathrm{ECA}}$ とし、
$U\in\mathfrak U_{\mathrm{ECA}}$ に対して公理的に取り扱う候補集合
$\operatorname{Ax}(U)$ を対応させる。

大域的公理候補集合は
\[
\Ax=\bigcup_{U\in\mathfrak U_{\mathrm{ECA}}}\operatorname{Ax}(U)
\]
とする。

公理選択空間は
\[
I\subseteq\mathcal{P}(\Ax)
\]
であり、Solver は
\[
s=(\mathcal{A}_s,\mathcal{R}_s),
\qquad
\Solver_I=\{s\mid\mathcal{A}_s\in I\}
\]
とする。

また、
\[
\mathcal{Z}=
\left\{
s\in\Solver_I\mid
\mathcal{T}(s)\cong_{\mathrm{Th}}\mathrm{ZFC}
\right\}.
\]

\section{$\mathcal{R}^*$ 選択閉包条件}
$\mathcal{R}^*\subseteq\operatorname{Ax}(U)\subseteq\Ax$ が構成されたと仮定し、
ZFC を表現する基礎公理集合を
\[
\mathcal{A}_{\mathrm{ZFC}}\subseteq\Ax
\]
とする。

\begin{definition}
\textbf{$\mathcal{R}^*$ 選択閉包条件}を
\[
\boxed{
\forall A\subseteq\mathcal{R}^*,
\qquad
\mathcal{A}_{\mathrm{ZFC}}\cup A\in I
}
\]
と定義する。
\end{definition}

これは
$\mathcal{P}(\mathcal{R}^*)\subseteq I$
を要求するものではない。必要なのは
\[
\Gamma:\mathcal{P}(\mathcal{R}^*)\to I,
\qquad
\Gamma(A)=\mathcal{A}_{\mathrm{ZFC}}\cup A
\]
という写像が定義できることである。

\section{ECA の現行定義から導出できるか}
\[
I\subseteq\mathcal{P}(\Ax)
\]
は、$I$ の各要素が $\Ax$ の部分集合であることしか述べていない。
したがって、
\[
\forall A\subseteq\mathcal{R}^*,
\quad
\mathcal{A}_{\mathrm{ZFC}}\cup A\in I
\]
は現行定義からは導出できない。

つまり、
\[
\mathcal{P}(\mathcal{R}^*)\longrightarrow I
\]
という全部分集合の選択可能性は現時点では未確定である。

\section{最小限の追加条件}
大域条件
\[
I=\mathcal{P}(\Ax)
\]
を仮定する必要はない。対象となる $\mathcal{R}^*$ に対してのみ
\[
\forall A\subseteq\mathcal{R}^*,
\quad
\mathcal{A}_{\mathrm{ZFC}}\cup A\in I
\]
を要求すればよい。

さらに
\[
\mathcal{R}^*\cap\mathcal{A}_{\mathrm{ZFC}}=\varnothing
\]
なら、
\[
A\neq B\Longrightarrow
\mathcal{A}_{\mathrm{ZFC}}\cup A\neq
\mathcal{A}_{\mathrm{ZFC}}\cup B
\]
である。

\begin{proposition}
\[
|\mathcal{R}^*|=\aleph_0,
\quad
\mathcal{R}^*\cap\mathcal{A}_{\mathrm{ZFC}}=\varnothing
\]
および $\mathcal{R}^*$ 選択閉包条件を仮定すると、
\[
\boxed{|I|\ge2^{\aleph_0}}
\]
が成り立つ。
\end{proposition}

\begin{proof}
\[
\Gamma(A)=\mathcal{A}_{\mathrm{ZFC}}\cup A
\]
とする。上記の条件により $\Gamma$ は単射である。
また
\[
|\mathcal{P}(\mathcal{R}^*)|
=2^{\aleph_0}.
\]
したがって
\[
2^{\aleph_0}
=
|\mathcal{P}(\mathcal{R}^*)|
\le |I|.
\]
\end{proof}

\section{反例}
例えば
\[
I=\{\mathcal{A}_{\mathrm{ZFC}}\}
\]
は
\[
I\subseteq\mathcal{P}(\Ax)
\]
に矛盾しない。しかし一般には
\[
\mathcal{A}_{\mathrm{ZFC}}\cup A\in I
\]
は $A=\varnothing$ の場合しか成立しない。

したがって
\[
I\subseteq\mathcal{P}(\Ax)
\]
だけでは
\[
|I|\ge2^{\aleph_0}
\]
も
\[
\mathcal{P}(\mathcal{R}^*)\hookrightarrow I
\]
も保証されない。

これは、$\operatorname{Ax}(U)$ に多数の候補が存在することと、それらの任意の組合せが $I$ で許容されることが別条件であることを示す。

\section{有限部分集合による検証}
\[
A_n=\{r_0,\ldots,r_{n-1}\}
\]
について
\[
\forall n\in\mathbb{N},\quad
\mathcal{A}_{\mathrm{ZFC}}\cup A_n\in I
\]
という有限選択可能性を考えることができる。

ただし、
\[
\text{有限部分集合に対する閉包}
\not\Rightarrow
\text{全部分集合に対する閉包}
\]
である。したがって、有限的な閉包性だけから全部分集合に対する条件を導くことはできない。

\section{局所条件と大域条件}
局所条件は
\[
\forall A\subseteq\mathcal{R}^*,
\quad
\mathcal{A}_{\mathrm{ZFC}}\cup A\in I
\]
である。

一方、大域条件は
\[
I=\mathcal{P}(\Ax)
\]
である。

大域条件は局所条件を含意する場合があるが、局所条件から大域条件は導けないため、今回の目的ではまず局所条件を検証する方が適切である。

\section{Stage 2 と Stage 3 の境界}
本稿は Stage 2、
\[
\mathcal{A}_{\mathrm{ZFC}}\cup A\in I
\]
までを扱う。

ここから直ちに
\[
(\mathcal{A}_{\mathrm{ZFC}}\cup A,\mathcal{R}_0)\in\mathcal{Z}
\]
は得られない。

Stage 3 ではさらに
\[
\mathcal{T}(\mathcal{A}_{\mathrm{ZFC}}\cup A,\mathcal{R}_0)
\cong_{\mathrm{Th}}\mathrm{ZFC}
\]
を要求する必要がある。

したがって現段階では
\[
\mathcal{R}^*\overset{\mathrm{Stage\ 2}}{\longrightarrow}I
\]
であり、
\[
I\overset{\mathrm{Stage\ 3}}{\longrightarrow}\mathcal{Z}
\]
は別の検証となる。

\section{連続体濃度に関する位置づけ}
本稿で得られる条件付き結果は
\[
|\mathcal{R}^*|=\aleph_0
\land
\text{$\mathcal{R}^*$ 選択閉包条件}
\Longrightarrow
|I|\ge2^{\aleph_0}
\]
である。

しかし、
\[
|I|\ge2^{\aleph_0}
\]
から
\[
|\mathcal{Z}|\ge2^{\aleph_0}
\]
を導くことはできない。 $\mathcal{Z}$ はさらに
\[
\mathcal{T}(s)\cong_{\mathrm{Th}}\mathrm{ZFC}
\]
という条件で切り出されるからである。

よって
\[
\mathcal{R}^*
\to
\mathcal{P}(\mathcal{R}^*)
\to
I
\to
\Solver_I
\to
\mathcal{Z}
\]
の各段階を別々に検証する必要がある。

\section{現段階での結論}
ECA の現在の定義から確定しているのは
\[
I\subseteq\mathcal{P}(\Ax)
\]
である。

これだけでは
\[
\mathcal{P}(\mathcal{R}^*)\hookrightarrow I
\]
は導出できない。

一方、
\[
\mathcal{R}^*\cap\mathcal{A}_{\mathrm{ZFC}}=\varnothing
\]
と
\[
\forall A\subseteq\mathcal{R}^*,
\quad
\mathcal{A}_{\mathrm{ZFC}}\cup A\in I
\]
を追加条件として採用すれば、
\[
\mathcal{P}(\mathcal{R}^*)\hookrightarrow I
\]
および
\[
|I|\ge2^{\aleph_0}
\]
を導出できる。

次は、単に $I$ を大きく仮定することではなく、
\[
I,\quad B(i),\quad\operatorname{Ax}(U),\quad\mathcal{T}
\]
の既存構造から、この局所的な選択閉包性を導出できるかどうかを考える。

% ECA_21.tex
\section{ECA_21 $I$ の選択閉包性と $B(i)$ の役割}

\section{目的}

ECA_20 では、可算無限公理候補族
\[
\mathcal{R}^*=\{r_n\mid n\in\mathbb{N}\}
\]
に対して
\[
\forall A\subseteq\mathcal{R}^*,
\qquad
\mathcal{A}_{\mathrm{ZFC}}\cup A\in I
\]
という局所的な選択閉包条件を考えた。

本節では、この条件が ECA から
\[
I\subseteq\mathcal{P}(\Ax),
\qquad
B(i)\subseteq i\subseteq\Ax
\]
自然に導出される可能性についてを検証する。

特に、
\[
B(i)
\]
が単なる部分集合として定義されているだけなのか、それとも公理選択集合 $i$ の構成規則に関する条件を与えるものなのかを区別する。

\section{現行の基本関係}

公理候補集合を
\[
\Ax
\]
とし、公理選択空間を
\[
I\subseteq\mathcal{P}(\Ax)
\]
とする。

各
\[
i\in I
\]
について
\[
B(i)\subseteq i\subseteq\Ax
\]
を満たす。

また、Solver
\[
s=(\mathcal{A}_s,\mathcal{R}_s)
\]
について
\[
\mathcal{A}_s\in I
\]
を要求する。

したがって、
\[
B(i)
\]
は少なくとも形式上は、選択された公理集合 $i$ の内部にある部分集合である。

\section{$B(i)$ から全部分集合閉包を導けるか}

現在の関係
\[
B(i)\subseteq i
\]
だけから、
\[
A\subseteq i
\Longrightarrow A\in I
\]
を導くことはできない。

そもそも $B(i)$ の定義は
\[
B(i)\subseteq i
\]
であり、
\[
B(i)=\mathcal{P}(i)
\]
でも
\[
B(i)=\{A\subseteq i\mid A\in I\}
\]
でもない。

したがって、現行の定義だけでは $B(i)$ によって $I$ の閉包性を表現することはできない。

\begin{proposition}
\[
I\subseteq\mathcal{P}(\Ax),
\qquad
B(i)\subseteq i
\]
だけでは、
\[
\forall A\subseteq\mathcal{R}^*,
\quad
\mathcal{A}_{\mathrm{ZFC}}\cup A\in I
\]
は導出できない。
\end{proposition}

\begin{proof}
例えば
\[
I=\{\mathcal{A}_{\mathrm{ZFC}}\}
\]
とし、
\[
B(\mathcal{A}_{\mathrm{ZFC}})
=
\varnothing
\]
とすれば、
\[
B(i)\subseteq i\subseteq\Ax
\]
は成立する。

一方、$\mathcal{R}^*\neq\varnothing$ なら、一般に
\[
\mathcal{A}_{\mathrm{ZFC}}\cup A
\notin I
\]
となる $A\subseteq\mathcal{R}^*$ が存在する。

よって、現在の包含関係だけでは選択閉包性は導けない。
\end{proof}

\section{$B(i)$ を選択可能部分族として読む場合}

別の定義として、
\[
B(i)
=
\{A\subseteq i\mid A\text{ が許容される}\}
\]
のように、$B(i)$ を $i$ 内部の選択可能部分集合族として定義することは考えられる。

この場合には
\[
B(i)\subseteq\mathcal{P}(i)
\]
となる。

さらに
\[
B(i)=\mathcal{P}(i)
\]
を仮定すれば、$i$ の任意の部分集合が $B(i)$ に属する。

しかし、これは現在の
\[
B(i)\subseteq i
\]
とは別の定義である。

したがって、この解釈を ECA に導入する場合は、既存の $B(i)$ の意味を変更することになるため、現時点では導入せず、別案として扱う。

\section{$I$ の局所閉包性を直接定義する場合}
$B(i)$ を変更せずに、$I$ の閉包性を直接記述する方法がある。

対象となる
\[
\mathcal{R}^*\subseteq\Ax
\]
と
\[
\mathcal{A}_{\mathrm{ZFC}}\subseteq\Ax
\]
に対して、
\[
\boxed{
\forall A\subseteq\mathcal{R}^*,
\quad
\mathcal{A}_{\mathrm{ZFC}}\cup A\in I
}
\]
を直接条件として置く。

この方法では
\[
B(i)
\]
の意味を変更する必要がない。

また、
\[
I=\mathcal{P}(\Ax)
\]
という大域的な条件も不要である。

したがって、現行構造をできるだけ保存するという観点では、この記述が最も直接的である。

\section{$B(i)$ と局所閉包条件の関係}

仮に
\[
i_A=\mathcal{A}_{\mathrm{ZFC}}\cup A
\]
と置く。

ECA_20 の局所閉包条件は
\[
\forall A\subseteq\mathcal{R}^*,
\quad i_A\in I
\]
である。

一方、
\[
B(i_A)\subseteq i_A
\]
は各 $i_A$ が選択された後に成立する包含関係である。

したがって、論理的には

\[
\boxed{
i_A\in I
\quad\text{と}\quad
B(i_A)\subseteq i_A
}
\]

は別の情報である。

前者は「その公理集合が選択空間 $I$ に属する」ことを表し、後者は「その公理集合に対応する $B$ が $i_A$ の部分集合である」ことを表す。

このため、現在の定義では $B(i)$ を用いて Stage 2 の選択閉包性を導く必要はない。

\section{有限閉包と全部分集合閉包}

$B(i)$ に何らかの閉包性を追加する場合にも、有限の場合と全部分集合の場合を区別する必要がある。

例えば
\[
A,B\subseteq\mathcal{R}^*
\]
について
\[
i_A,i_B\in I
\]
から
\[
i_{A\cup B}\in I
\]
を導く二項閉包性を考えることができる。

しかし、このような有限的な閉包性だけでは
\[
\forall A\subseteq\mathcal{R}^*,
\quad i_A\in I
\]
を直ちに導けない。

特に可算無限集合
\[
A=\{r_n\mid n\in J\},
\qquad J\subseteq\mathbb{N}
\]
については、有限近似
\[
A_n=A\cap\{r_0,\ldots,r_{n-1}\}
\]
から
\[
i_A\in I
\]
を得るための極限・可算閉包に相当する条件が別途必要になる。

\section{選択閉包を ECA の構造条件とする場合}

ECA の基本構造を変更することなく、選択空間 $I$ に対する追加条件として
\[
\mathsf{Closure}(\mathcal{R}^*,\mathcal{A}_{\mathrm{ZFC}})
\]
を導入し、
\[
\mathsf{Closure}(\mathcal{R}^*,\mathcal{A}_{\mathrm{ZFC}})
\iff
\forall A\subseteq\mathcal{R}^*,
\quad
\mathcal{A}_{\mathrm{ZFC}}\cup A\in I
\]
と定義することができる。

この場合、
\[
\mathsf{Closure}
\]
は $B(i)$ の定義ではなく、$I$ に対する追加の構造条件となる。

この分離には、

\[
\operatorname{Ax}(U)
\quad\longrightarrow\quad
\mathcal{R}^*
\]

と

\[
\mathcal{R}^*
\quad\longrightarrow\quad
I
\]

を別々に検証できるという利点がある。

\section{連続体濃度への影響}

\[
|\mathcal{R}^*|=\aleph_0
\]
および
\[
\mathsf{Closure}(\mathcal{R}^*,\mathcal{A}_{\mathrm{ZFC}})
\]
に加えて
\[
\mathcal{R}^*\cap\mathcal{A}_{\mathrm{ZFC}}=\varnothing
\]
を仮定すると、
\[
\Gamma:
\mathcal{P}(\mathcal{R}^*)\hookrightarrow I,
\qquad
\Gamma(A)=\mathcal{A}_{\mathrm{ZFC}}\cup A
\]
が得られる。

したがって
\[
|I|\ge2^{\aleph_0}.
\]

しかし、これは依然として
\[
|\mathcal Z|\ge2^{\aleph_0}
\]
を意味しない。

Stage 3 において
\[
\mathcal T(s_A)\cong_{\mathrm{Th}}\mathrm{ZFC}
\]
をすべての $A\subseteq\mathcal{R}^*$ について確認する必要がある。

\section{現段階で採用すべき構造}

現時点では、
\[
B(i)\subseteq i
\]
という既存の定義を変更せず、

\[
\boxed{
\forall A\subseteq\mathcal{R}^*,
\quad
\mathcal{A}_{\mathrm{ZFC}}\cup A\in I
}
\]

を独立した Stage 2 の検証条件として扱うのが最も好ましい。

これにより、

\[
\operatorname{Ax}(U)
\to
\mathcal{R}^*
\]

は ECA_19、

\[
\mathcal{R}^*
\to
I
\]

はこの文書、そして

\[
I
\to
\mathcal Z
\]

は Stage 3 の内容として分離できる。

\section{結論}

現在の ECA における
\[
I\subseteq\mathcal{P}(\Ax),
\qquad
B(i)\subseteq i
\]
から、$\mathcal{R}^*$ の全部分集合に対する選択閉包性を導くことはできない。

特に、
\[
B(i)\subseteq i
\]
は、それ自体では $I$ の閉包性を規定していない。

したがって、現段階では
\[
B(i)
\]
の意味を変更して閉包性を担わせるより、

\[
\boxed{
\forall A\subseteq\mathcal{R}^*,
\quad
\mathcal{A}_{\mathrm{ZFC}}\cup A\in I
}
\]
を独立した選択構造上の条件として検証する方が良いと判断できる。

% ECA_22.tex
\section{ECA_22 公理選択空間 $I$ の定義と選択閉包性}

\section{目的}

ECA_21 では、
\[
I\subseteq\mathcal{P}(\Ax),
\qquad
B(i)\subseteq i\subseteq\Ax
\]
という既存構造だけから、
\[
\forall A\subseteq\mathcal{R}^*,
\qquad
\mathcal{A}_{\mathrm{ZFC}}\cup A\in I
\]
を導出することはできないことを確認した。

そこで本稿では、一旦 $I$ の定義そのものから考え「ECA において許容される公理選択とは何か」を検証する。
ここでの目的は現在の ECA における $I$ の構造から、どの程度の選択閉包性が自然に得られるかを明確にすることである。

\section{$I$ の基本的位置}

ECA では $\Ax$ を公理および公理スキームとして扱う候補の全体とし
\[
I\subseteq\mathcal{P}(\Ax)
\]
とする。

したがって、$i\in I$ は $\Ax$ の部分集合であり、ECA において公理選択として扱われる集合である。
ここで重要なのは、$\Ax$ と $I$ が異なる対象であることである。

$a\in\Ax$ は公理候補そのものを表す一方、$i\in I$ は公理候補を組み合わせた選択集合を表す。

したがって、これは
\[
\Ax=\text{候補の宇宙},
\qquad
I=\text{許容される選択集合の空間}
\]
という役割分担になる。

\section{部分集合としての $I$}

\[
I\subseteq\mathcal{P}(\Ax)
\]
は、$I$ が $\Ax$ の全部分集合を必ず含むという意味ではない。

例えば、
\[
I=\{\varnothing,\mathcal{A}_0\}
\]
も形式的には
\[
I\subseteq\mathcal{P}(\Ax)
\]
を満たし得る。

したがって、
\[
\forall A\subseteq\Ax,\quad A\in I
\]
を要求する
\[
I=\mathcal{P}(\Ax)
\]
とは明確に異なる。

この区別は、$\mathcal{R}^*$ の全部分集合を選択可能にする問題に直接関係する。

\section{許容性を表す条件}

$I$ を単なる任意の部分集合として扱うのではなく、ある許容性条件によって定めることを考える。

抽象的に
\[
\operatorname{Adm}(i)
\]
を「$i$ が ECA において許容される公理選択である」という条件とする。

このとき、
\[
\boxed{
I=\{i\subseteq\Ax\mid\operatorname{Adm}(i)\}
}
\]
と表現できる。

この表現を採用した場合、ECA の問題は
\[
\operatorname{Adm}(i)
\]
をどのように定義するかへ移る。

ただし、これは $I$ を形式化するための一般形として扱う。

\section{$B(i)$ との関係}

現在の
\[
B(i)\subseteq i
\]
は、$B(i)$ が $i$ の内部にあることを表す。

一方、
\[
\operatorname{Adm}(i)
\]
は $i$ 自体が選択空間 $I$ に属するための条件である。

したがって、
\[
B(i)\subseteq i
\]
と
\[
\operatorname{Adm}(i)
\]
は異なる役割を持つ。

概念的には、

\[
\operatorname{Adm}(i)
\longrightarrow
i\in I
\]

と

\[
B(i)\subseteq i
\]

を分けて扱うことができる。

この分離によって $I$ の選択規則を考えることができる。

\section{$\mathcal{R}^*$ に対する許容性}

ECA_19 により、
\[
\mathcal{R}^*
\subseteq\operatorname{Ax}(U)
\subseteq\Ax
\]
が得られたと仮定する。

ZFC を表現する基礎集合を
\[
\mathcal{A}_{\mathrm{ZFC}}\subseteq\Ax
\]
とする。

このとき
\[
i_A=\mathcal{A}_{\mathrm{ZFC}}\cup A,
\qquad
A\subseteq\mathcal{R}^*
\]
を考える。

Stage 2 の条件は
\[
\boxed{
\forall A\subseteq\mathcal{R}^*,
\quad
\operatorname{Adm}(i_A)
}
\]
と書き換えることができる。

すなわち、
\[
\forall A\subseteq\mathcal{R}^*,
\quad
i_A\in I
\]
を保証するためには、
\[
\operatorname{Adm}
\]
がこの族に対してどのような性質を持つかを検討すればよい。

\section{許容性の候補となる性質}

ここでは、検証対象としていくつかの性質を区別する。

\subsection{単純な包含閉包}

\[
i\in I,\quad A\subseteq i
\Longrightarrow
A\in I.
\]

これは $I$ が部分集合について閉じていることを意味する。

しかし、
\[
\mathcal{A}_{\mathrm{ZFC}}\cup A
\]
は $\mathcal{A}_{\mathrm{ZFC}}$ の部分集合ではないため、この性質だけでは
\[
\mathcal{A}_{\mathrm{ZFC}}\cup A\in I
\]
を導けない。

\subsection{有限和閉包}

\[
i,j\in I
\Longrightarrow
i\cup j\in I.
\]

これは有限個の許容された選択集合を組み合わせる条件である。

ただし、これだけでは
\[
\bigcup_{n\in\mathbb{N}}i_n\in I
\]
を導けない。

\subsection{可算和閉包}

\[
(i_n)_{n\in\mathbb{N}}\subseteq I
\Longrightarrow
\bigcup_{n\in\mathbb{N}}i_n\in I.
\]

これは有限和閉包より強い条件である。

しかし、これだけでは「各 $A\subseteq\mathcal{R}^*$ に対する $\mathcal{A}_{\mathrm{ZFC}}\cup A$」が $I$ に属することを直ちに意味しない。

\subsection{局所的な $\mathcal{R}^*$ 閉包}

最も直接的なのは、
\[
\boxed{
\forall A\subseteq\mathcal{R}^*,
\quad
\operatorname{Adm}
(\mathcal{A}_{\mathrm{ZFC}}\cup A)
}
\]
を要求することである。

これは ECA_20からECA_21 で扱った条件そのものであり、必要な範囲に限定された条件である。

\section{閉包性だけでは十分でない場合}

$I$ に閉包性を与えても、
\[
\mathcal{A}_{\mathrm{ZFC}}\in I
\]
が成立しなければ、そこから
\[
\mathcal{A}_{\mathrm{ZFC}}\cup A\in I
\]
を構成することはできない。

したがって、例えば有限和閉包を採用する場合には

\[
\mathcal{A}_{\mathrm{ZFC}}\in I
\]

と $mathbb{N}$ ($\{r\}$) を組み込むための適切な許容性

など、出発点となる条件も必要になる。

このため、
\[
\text{候補の存在}
\]
と
\[
\text{候補を含む選択集合の許容性}
\]
は分離して検証する必要がある。

\section{必要条件と十分条件の分離}

Stage 2 に関して、次の条件を区別する。

\subsection*{候補存在}
\[
\exists U\in\mathfrak U_{\mathrm{ECA}},
\quad
\exists\mathcal{R}^*
\subseteq\operatorname{Ax}(U),
\quad
|\mathcal{R}^*|=\aleph_0.
\]

これは ECA_19 の対象である。

\subsection*{基礎選択}
\[
\mathcal{A}_{\mathrm{ZFC}}\in I.
\]

\subsection*{局所閉包}
\[
\forall A\subseteq\mathcal{R}^*,
\quad
\mathcal{A}_{\mathrm{ZFC}}\cup A\in I.
\]

この三つは同一の条件ではない。

特に、
\[
|\operatorname{Ax}(U)|\ge\aleph_0
\]
から
\[
|I|\ge2^{\aleph_0}
\]
を直接導くことはできない。

\section{連続体濃度との関係}

\[
|\mathcal{R}^*|=\aleph_0
\]
かつ
\[
\mathcal{R}^*\cap\mathcal{A}_{\mathrm{ZFC}}=\varnothing
\]
であり、局所閉包条件が成立すれば、

\[
\Gamma:
\mathcal{P}(\mathcal{R}^*)\hookrightarrow I,
\qquad
\Gamma(A)=\mathcal{A}_{\mathrm{ZFC}}\cup A
\]
が得られる。

したがって、
\[
|I|\ge2^{\aleph_0}.
\]

しかし、ここでも
\[
|I|\ge2^{\aleph_0}
\not\Rightarrow
|\mathcal Z|\ge2^{\aleph_0}
\]
である。

\[
\mathcal Z
\]
については、さらに
\[
\mathcal T(s)\cong_{\mathrm{Th}}\mathrm{ZFC}
\]
を満たす Solver を構成する必要がある。

\section{現段階における可能性}

現行 ECA の
\[
I\subseteq\mathcal{P}(\Ax)
\]
という定義だけからは、$\mathcal{R}^*$ の全部分集合に対する選択可能性は導出できない可能性がある。

また、
\[
B(i)\subseteq i
\]
をそのための閉包性として読み替える必要もない。

したがって、現段階では

\[
\boxed{
I=\{i\subseteq\Ax\mid\operatorname{Adm}(i)\}
}
\]

という一般的な許容性表現を導入して検討するのが適切である。

その上で、必要最小限の条件として

\[
\boxed{
\forall A\subseteq\mathcal{R}^*,
\quad
\operatorname{Adm}
(\mathcal{A}_{\mathrm{ZFC}}\cup A)
}
\]

を検証対象とする。

この条件が ECA の既存構造から導出されるなら追加公理として扱う必要性は低くなり、導出できないなら、ECA において明示的な選択条件として採用するかどうかを改めて判断することになる。

\section{次段階}

次に必要なのは、Stage 3 の検証である。

すなわち、
\[
i_A=\mathcal{A}_{\mathrm{ZFC}}\cup A\in I
\]
から
\[
s_A=(i_A,\mathcal{R}_0)\in\mathcal Z
\]
を得るために、

\[
\mathcal T(i_A,\mathcal{R}_0)
\cong_{\mathrm{Th}}
\mathrm{ZFC}
\]
をすべての
\[
A\subseteq\mathcal{R}^*
\]
について保証できる条件を検討する。

ここでは、単に各 $r_n$ が個別に ZFC と矛盾しないという条件ではなく、任意の部分集合 $A$ を加えた場合にも同じ数学的構造表現を保てることが必要になる。

\section{ECA_23 \(\mathcal{R}^*\) の部分集合に対する ZFC 保存条件}


\section{目的}
ECA_22 までにより、可算無限公理候補族
\[
\mathcal{R}^*=\{r_n\mid n\in\mathbb{N}\}
\]
を考えたとき、Stage 2 は
\[
i_A=\mathcal{A}_{\mathrm{ZFC}}\cup A\in I
\qquad(A\subseteq\mathcal{R}^*)
\]
という選択可能性の問題として整理された。

本節では Stage 3、
\[
i_A\in I\longrightarrow s_A\in\mathcal{Z}
\]
を検証する。各 \(r_n\) が個別に ZFC と両立することだけでは不十分であり、任意の \(A\subseteq\mathcal{R}^*\) について
\[
\mathcal{T}(s_A)\cong_{\mathrm{Th}}\mathrm{ZFC}
\]
を保証する必要がある。

\section{Stage 3 の基本設定}
Stage 2 で
\[
i_A=\mathcal{A}_{\mathrm{ZFC}}\cup A
\]
が \(I\) に属するとする。固定した推論規則・構成制約を
\[
\mathcal{R}_0
\]
とし、
\[
s_A=(i_A,\mathcal{R}_0)
\]
と定義する。

数学的構造写像
\[
\mathcal{T}:\Solver_I\to\mathrm{Th}
\]
に対して、
\[
\mathcal{Z}=
\left\{
s\in\Solver_I\mid
\mathcal{T}(s)\cong_{\mathrm{Th}}\mathrm{ZFC}
\right\}
\]
であるから、Stage 3 の条件は
\[
\boxed{
\forall A\subseteq\mathcal{R}^*,
\quad
\mathcal{T}(i_A,\mathcal{R}_0)
\cong_{\mathrm{Th}}\mathrm{ZFC}
}
\]
となる。

\section{個別的保存と族全体の保存}
候補 \(r\in\mathcal{R}^*\) について
\[
\mathcal{T}(\mathcal{A}_{\mathrm{ZFC}}\cup\{r\},\mathcal{R}_0)
\cong_{\mathrm{Th}}\mathrm{ZFC}
\]
が成立しても、一般には
\[
\forall A\subseteq\mathcal{R}^*,
\quad
\mathcal{T}(\mathcal{A}_{\mathrm{ZFC}}\cup A,\mathcal{R}_0)
\cong_{\mathrm{Th}}\mathrm{ZFC}
\]
は導けない。

複数の候補を同時に追加することで新しい結論が生じる可能性があるためである。

\begin{definition}
\(\mathcal{R}^*\) が \(\mathcal{A}_{\mathrm{ZFC}}\) および \(\mathcal{R}_0\) に関して \textbf{ZFC 保存族}であるとは、
\[
\boxed{
\forall A\subseteq\mathcal{R}^*,
\quad
\mathcal{T}(\mathcal{A}_{\mathrm{ZFC}}\cup A,\mathcal{R}_0)
\cong_{\mathrm{Th}}\mathrm{ZFC}
}
\]
が成立することとする。
\end{definition}

\section{保存条件の強さ}
\(A=\varnothing\) も対象なので、
\[
\mathcal{T}(\mathcal{A}_{\mathrm{ZFC}},\mathcal{R}_0)
\cong_{\mathrm{Th}}\mathrm{ZFC}
\]
が必要である。また \(A=\mathcal{R}^*\) も対象なので、
\[
\mathcal{T}(\mathcal{A}_{\mathrm{ZFC}}\cup\mathcal{R}^*,\mathcal{R}_0)
\cong_{\mathrm{Th}}\mathrm{ZFC}
\]
も必要である。さらに有限部分集合だけでなく、可算無限部分集合を含むすべての \(A\subseteq\mathcal{R}^*\) が対象になる。

\section{数学的構造同型と無矛盾性の区別}
\[
\mathcal{T}(\mathcal{A}_{\mathrm{ZFC}}\cup A,\mathcal{R}_0)
\cong_{\mathrm{Th}}\mathrm{ZFC}
\]
と単なる無矛盾性条件を同一視してはならない。

無矛盾であることだけでは、本節の \(\mathcal{Z}\) の定義に必要な数学的構造同型性を保証しない。

\section{モデルによる条件との区別}
数学的構造とモデルを区別する。
\[
M\models\mathcal{T}(s_A)
\]
はモデルについての条件であり、
\[
\mathcal{T}(s_A)\cong_{\mathrm{Th}}\mathrm{ZFC}
\]
はSolverの数学的構造についての条件である。

したがって、
\[
\exists M\,
(M\models\mathcal{T}(s_A)\land M\models\mathrm{CH})
\]
が成立しても、それだけから \(s_A\in\mathcal{Z}\) は導かれない。

\section{Solver の推論規則を固定する意味}
\[
s_A=(i_A,\mathcal{R}_0)
\]
として \(\mathcal{R}_0\) を固定するのは、今回の検証における変化を公理選択部分に限定するためである。

もし
\[
s_A=(i_A,\mathcal{R}_A)
\]
として推論規則・構成制約まで変更すると、公理選択と Solver 構造の差を同時に扱うことになる。

これは ECA の一般定理ではなく、Stage 3 を明確に検証するための設定である。

\section{Stage 2 と Stage 3}
Stage 2:
\[
\forall A\subseteq\mathcal{R}^*,
\quad
\mathcal{A}_{\mathrm{ZFC}}\cup A\in I.
\]

Stage 3:
\[
\forall A\subseteq\mathcal{R}^*,
\quad
\mathcal{T}(\mathcal{A}_{\mathrm{ZFC}}\cup A,\mathcal{R}_0)
\cong_{\mathrm{Th}}\mathrm{ZFC}.
\]

前者により \(s_A\in\Solver_I\) が成立し、後者により \(s_A\in\mathcal{Z}\) が成立する。したがって、
\[
\boxed{\forall A\subseteq\mathcal{R}^*,\quad s_A\in\mathcal{Z}}
\]
を得るための二段階条件として整理できる。

\section{連続体濃度への接続}
さらに
\[
|\mathcal{R}^*|=\aleph_0,
\qquad
\mathcal{R}^*\cap\mathcal{A}_{\mathrm{ZFC}}=\varnothing
\]
を仮定する。

\[
\Phi:\mathcal{P}(\mathcal{R}^*)\to\mathcal{Z},
\qquad
\Phi(A)=
(\mathcal{A}_{\mathrm{ZFC}}\cup A,\mathcal{R}_0)
\]
と定義する。

\(A\neq B\) なら
\[
\mathcal{A}_{\mathrm{ZFC}}\cup A
\neq
\mathcal{A}_{\mathrm{ZFC}}\cup B
\]
なので \(\Phi\) は単射である。従って、上記条件がすべて成立すれば
\[
\boxed{
|\mathcal{Z}|
\ge
|\mathcal{P}(\mathcal{R}^*)|
=
2^{\aleph_0}
}
\]
となる。

これは条件付き結果であり、ECA の現行定義から自動的に得られる結果ではない。

\section{CH との分離}
ここで得られる
\[
|\mathcal{Z}|\ge2^{\aleph_0}
\]
は
\[
2^{\aleph_0}=\aleph_1
\]
を主張しない。したがって、\(\mathcal{Z}\) の濃度問題と CH は別の命題として扱う。

\section{現時点で未保証の事項}
現行 ECA の定義だけから、次の全てが自動的に得られるとは言えない。
\[
|\mathcal{R}^*|=\aleph_0,
\]
\[
\forall A\subseteq\mathcal{R}^*,
\quad
\mathcal{A}_{\mathrm{ZFC}}\cup A\in I,
\]
\[
\forall A\subseteq\mathcal{R}^*,
\quad
\mathcal{T}(\mathcal{A}_{\mathrm{ZFC}}\cup A,\mathcal{R}_0)
\cong_{\mathrm{Th}}\mathrm{ZFC}.
\]

ECA_19 は第一段階、ECA_20～ECA_22 は第二段階、本節は第三段階を担当する。

\section{今後の検証点}
Stage 3 を具体化するには、
\begin{enumerate}
\item \(r\in\mathcal{R}^*\) が ZFC 保存候補となるための条件
\item 個々の候補の保存性から族全体の保存性を導く条件
\item $\mathcal{T}(\mathcal{A}_{\mathrm{ZFC}}\cup A,\mathcal{R}_0)$ と $\mathcal{T}(\mathcal{A}_{\mathrm{ZFC}},\mathcal{R}_0)$ の同型を構成する方法
\item Solver の数学的構造表現と通常の数学的構造における保存性との関係
\item これらが ECA から導出できるか、または追加条件として明示すべきか
\end{enumerate}

\section{結論}
Stage 3 の本質は、各 \(r_n\) が単独で ZFC と両立することではない。

必要なのは
\[
\boxed{
\forall A\subseteq\mathcal{R}^*,
\quad
\mathcal{T}(\mathcal{A}_{\mathrm{ZFC}}\cup A,\mathcal{R}_0)
\cong_{\mathrm{Th}}\mathrm{ZFC}
}
\]
という族全体の ZFC 保存性である。

この条件と Stage 2 を合わせることで、
\[
\mathcal{P}(\mathcal{R}^*)\hookrightarrow\mathcal{Z}
\]
という連続体濃度の下界へ条件付きで接続できる。




\section{ECA_24 ZFC 保存候補の具体化と冗長な公理による部分集合族の保存条件}

\section{目的}

本稿では、ECA_23 で導入した
「$\mathcal{R}^*$ の部分集合を追加しても ZFC と同型な数学的構造が得られる」という条件をより具体的な論理条件へ分解する。

ここで扱うのは、Solver 集合版 ECA の構造から $\mathcal{R}^*$ や連続体濃度が既に得られているという主張ではない。
目的は、既に導入された
\[
\mathcal{A}_{\mathrm{ZFC}},\quad
\mathcal{R}^*,\quad
\mathcal{R}_0,\quad
i_A=\mathcal{A}_{\mathrm{ZFC}}\cup A,\quad
s_A=(i_A,\mathcal{R}_0)
\]
に対して、Stage 3 の
\[
s_A\in\mathcal{Z}
\]
を成立させるための条件を具体化することである。

Solver は ECA の構築過程から形成された ECA の構造であり、本節ではそれを前提として議論する。

\section{Solverの数学的構造の基礎}

基礎となる Solver を
\[
s_0=(\mathcal{A}_{\mathrm{ZFC}},\mathcal{R}_0)
\]
とし、その数学的構造を
\[
T_0:=\mathcal{T}(s_0)
=\mathcal{T}(\mathcal{A}_{\mathrm{ZFC}},\mathcal{R}_0)
\]
と置く。

ECA_23 における候補 Solver は
\[
s_A=(\mathcal{A}_{\mathrm{ZFC}}\cup A,\mathcal{R}_0),
\qquad
A\subseteq\mathcal{R}^*.
\]

したがって、追加される候補 $A$ が基礎となるSolverの数学的構造 $T_0$ に対してどのような関係を持つかを調べれば、Stage 3 の保存条件を具体化できる。

\section{単一候補の冗長性}
$r\in\mathcal{R}^*$ を一つ固定する。
$r$ が $T_0$ と同一の形式言語上の文として扱える場合、$T_0\vdash r$ を考えることができる。
これは、$r$ が $T_0$ において既に証明可能であることを意味する。
したがって、同一の言語・固定された推論規則の通常の公理追加として扱う限り $r$ を公理として追加しても数学的構造の定理集合は変化しない。

ただし、これには次の前提がある。

\begin{enumerate}
\item $r$ が $T_0$ と同じ形式言語上で表現されること
\item $r$ が一つの文として扱えること
\item $\mathcal{T}$ が公理集合への文の追加を通常の公理追加として解釈すること
\item 推論規則 $\mathcal{R}_0$ が変更されないこと
\end{enumerate}

したがって、$T_0\vdash r$ は ECA の現行定義から自動的に得られる条件ではなく、ZFC 保存を実現するための具体的な十分条件の候補である。

\section{部分集合全体への拡張}

単一の $r$ について冗長性が成立するだけでは、任意の $A\subseteq\mathcal{R}^*$ に対する保存を記述するには不十分である。

そこで、
\[
\boxed{\forall r\in\mathcal{R}^*,\quad T_0\vdash r}
\]
を考える。

任意の
\[
A\subseteq\mathcal{R}^*
\]
を取ると、各 $r\in A$ について
\[
T_0\vdash r
\]
が成立する。

同一の言語・固定した推論規則による通常の公理追加として扱う限り、
\[
\operatorname{Th}(T_0\cup A)
=
\operatorname{Th}(T_0)
\]
となる。

ここで $\operatorname{Th}(T)$ は、数学的構造 $T$ から導出される定理全体を表す。

したがって、
\[
\mathcal{T}(\mathcal{A}_{\mathrm{ZFC}}\cup A,\mathcal{R}_0)
\]
は基礎となるSolverの数学的構造 $T_0$ と同じ定理集合を持つ。

\section{ZFC 同型との関係}
ECA_23 で定義した
\[
\mathcal{Z}=
\left\{
s\in\Solver_I
\mid
\mathcal{T}(s)\cong_{\mathrm{Th}}\mathrm{ZFC}
\right\}
\]
に対しては、基礎となるSolverの数学的構造 $T_0$ 自体について $T_0\cong_{\mathrm{Th}}\mathrm{ZFC}$ を確認する必要がある。

そのうえで、同一言語・同一推論規則の下で
\[
\operatorname{Th}(T_0\cup A) = \operatorname{Th}(T_0)
\]
が成立するならば、
\[
\mathcal{T}(\mathcal{A}_{\mathrm{ZFC}}\cup A,\mathcal{R}_0) \cong_{\mathrm{Th}} \mathrm{ZFC}
\]
を得るための具体的な経路が得られる。

したがって、次の条件をまとめることができる。
\[
\boxed{
\begin{aligned}
&T_0\cong_{\mathrm{Th}}\mathrm{ZFC},\\
&\forall r\in\mathcal{R}^*,\quad T_0\vdash r,\\
&\forall A\subseteq\mathcal{R}^*,\quad
\mathcal{A}_{\mathrm{ZFC}}\cup A\in I.
\end{aligned}}
\]

このとき、上記の前提の下で
\[
\forall A\subseteq\mathcal{R}^*,\quad s_A\in\mathcal{Z}
\]
を得る。

\section{保存概念の区別}
ZFC 保存を論じる際には、異なる性質を同一視してはならない。

\subsection{既に証明可能であること}
$T_0\vdash r$ は、$r$ が基礎となるSolverの数学的構造の定理であることを表す。

この条件は、同一の言語・固定された推論規則の通常の公理追加に対して、$r$ の追加が冗長であることを示すために利用できる。

\subsection{整合性}

\[
\operatorname{Con}(T_0+\{r\})
\]
は、$T_0+r$ が無矛盾であることを表す。

しかし、
\[
\operatorname{Con}(T_0+\{r\})
\]
だけから
\[
\operatorname{Th}(T_0+\{r\})
=
\operatorname{Th}(T_0)
\]
は得られない。

したがって、整合性は ZFC 保存そのものとは異なる。

\subsection{保存拡大}
拡張したSolverの数学的構造 $T_1$ が基礎となるSolverの数学的構造 $T_0$ に対して保存的であるとは、通常は基礎となる言語 (論理命題) における定理について
\[
T_1\vdash\varphi
\Longrightarrow
T_0\vdash\varphi
\]
が成立することを意味する。

これは
\[
\operatorname{Th}(T_1)=\operatorname{Th}(T_0)
\]
より弱い条件である場合がある。

したがって、保存的拡大であることだけから、ECA_23 の
\[
\mathcal{T}(s_A)\cong_{\mathrm{Th}}\mathrm{ZFC}
\]
を直ちに得られるとは限らない。

\subsection{Solverの数学的構造における同型}
ECA における
\[
T_1\cong_{\mathrm{Th}}T_0
\]
は、ECA_7 で導入したSolverの数学的構造における同型の定義に従って扱う。

そのため、「保存的である」ことと「ECA の意味でSolverの数学的構造において同型である」ことを同一視しないものとする。

\section{条件の関係}
同一の言語・固定された推論規則で、公理追加を通常の公理追加として扱うという仮定の下では $T_0\vdash r$ は $\operatorname{Th}(T_0+\{r\}) = \operatorname{Th}(T_0)$ を導くための強い条件として利用できる。

一方、$\operatorname{Con}(T_0+\{r\})$ だけでは同じ結論は得られない。

したがって、本節では
\[
\boxed{
\text{定理として既に導出可能}
\quad\Rightarrow\quad
\text{冗長な公理追加}
}
\]
を Stage 3 の具体的な十分条件として採用する。

\section{公理スキームの場合}
$r$ が単一の文ではなく公理スキームである場合、
単純な
\[
T_0\vdash r
\]
という記述はそのまま適用できない。

この場合には、スキームの各インスタンス、またはスキームをどのような対象として $\mathcal{T}$ が扱うかを明示する必要がある。

例えばスキーム $\Sigma$ の各インスタンスを
\[
\Sigma[\varphi]
\]
と書くなら、適切な範囲の $\varphi$ に対して
\[
T_0\vdash\Sigma[\varphi]
\]
が成立することを検討する必要がある。

ただし、これは ECA_19 で問題にした「公理スキームそのもの」と「そのインスタンス」を $\operatorname{Ax}(U)$ の要素としてどう表現するかという問題と直結する。

したがって、公理スキームを含む $\mathcal{R}^*$ については、単一文の場合の議論をそのまま一般化してはならない。

\section{Stage 2 と Stage 3 の統合}

Stage 2 は
\[
\forall A\subseteq\mathcal{R}^*,\quad
\mathcal{A}_{\mathrm{ZFC}}\cup A\in I.
\]

Stage 3 は
\[
\forall A\subseteq\mathcal{R}^*,\quad
\mathcal{T}(\mathcal{A}_{\mathrm{ZFC}}\cup A,\mathcal{R}_0)
\cong_{\mathrm{Th}}\mathrm{ZFC}.
\]

さらに本稿では、Stage 3 を保証する具体的な強条件の候補として
\[
T_0\cong_{\mathrm{Th}}\mathrm{ZFC}
\]
および
\[
\forall r\in\mathcal{R}^*,\quad T_0\vdash r
\]
を置いた。

このとき、同一の言語・固定された推論規則などの前提の下で
\[
\forall A\subseteq\mathcal{R}^*,\quad s_A\in\mathcal{Z}
\]
を得る。

さらに
\[
|\mathcal{R}^*|=\aleph_0
\]
かつ
\[
\mathcal{R}^*\cap\mathcal{A}_{\mathrm{ZFC}}=\varnothing
\]
ならば、
\[
\Phi:\mathcal{P}(\mathcal{R}^*)\to\mathcal{Z},
\qquad
\Phi(A)=s_A
\]
は単射となる。

したがって条件付きで
\[
|\mathcal{Z}|
\ge
|\mathcal{P}(\mathcal{R}^*)|
=
2^{\aleph_0}
\]
を得る。

\section{現時点で未確定な点}

本稿の条件は、現在の ECA から自動的に得られる定理としてではなく、$\mathcal{Z}$ 内で連続体濃度を実現するための具体的な検証として整理される。

特に未確定なのは次の点である。

\begin{enumerate}
\item $\operatorname{Ax}(U)$ から実際に可算無限な $\mathcal{R}^*$ を構成できるか
\item $\mathcal{A}_{\mathrm{ZFC}}\cup A\in I$ が全ての $A\subseteq\mathcal{R}^*$ について成立するか
\item 基礎 Solver $s_0$ が $T_0\cong_{\mathrm{Th}}\mathrm{ZFC}$ を満たすことをどのように確認するか
\item $\mathcal{R}^*$ の各候補を $T_0\vdash r$ という形で実際に構成できるか
\item 公理スキームを含む場合に、各インスタンスを $\mathcal{T}$ がどのように扱うか
\item 上記条件が ECA の既存定義から導出されるのか、それとも Solver 構造を具体化するための明示条件として定める必要があるのか
\end{enumerate}

\section{最後に}
$\operatorname{Ax}(U)$ の中から実際に $\mathcal{R}^*$ を構成できる候補を調べることで、ECA_19 で残された「可算無限公理候補族を $\operatorname{Ax}(U)$ から実際に取り出せるか」という問題と、ECA_24 で具体化した「その候補族を ZFC 保存族として扱えるか」を解決できる可能性がある。

特に、
\[
\text{ZFC の定理}
\longrightarrow
\text{ECA が扱う公理候補}
\longrightarrow
\mathcal{R}^*
\]
という関係性が、現在の $\operatorname{Ax}(U)$ の定義の中で成立するかを確認することができるかが重要な点として考えられる。































\end{document}
